\documentclass{ltjsbook}
\usepackage{docmute} 
\usepackage{../../myStyle} 

\begin{document}
\chapter{確率的プログラミング；状態空間モデル}
\label{x29_SSpace}
前回は体重の変化データを例に，データの平均点が変わるモデル，変化点を検出したり変化点を境に傾きが変わる線形モデルをみてきました。

同じように，今回も時系列的なデータを扱います。前回の最後に案内したように，時系列データに普通の回帰直線をそのまま当てはめることは適切ではありません。どのような注意点が必要なのかについて，少し考えてみましょう。

\section{時系列データの特徴}
時系列的なデータはどのようなときに得られるでしょうか。前回導入した時のように，誰か1人が日々の生活を記録し続けている，それも立派な時系列データです。
心理学には\keytermL{日誌法}{にっしほう}と呼ばれるデータ収集法があります。一言で行ってしまえば日記なのですが，日々多くの出来事や感情が到来するのを後から振り返って考えるのでは，正しく思い出せなかったり記憶が歪んでしまうこともあります。日々の状態を細かく記録し積み重ねることは，数字になっていなくても貴重なデータなのです。また最近は，\keytermL{経験サンプリング}{けいけんさんぷりんぐ}と呼ばれる調査法もあります。これは決まった間隔で質問紙がスマートフォンなどに送られてきて，定期的に心情を心理尺度に反映させていく，というものです。こうしたことができるようになったのには，私たちの周りに電子的なデバイスがたくさんあることが理由の1つです。さらに最近ではウェアラブル端末といって，身につける電子端末がありますね。Apple Watchなど携帯電話と同じような機能を持ったこれらのデバイスは，GPS機能がついていたり，万歩計による歩数の計測，血圧・脈拍のリアルタイムな計測が行われていたりします\footnote{こうしたデータは\keyterm{ライフログ}{らいふろぐ}{life-log}と呼ばれることもあります。}。これを蓄積したデータとすると，かなりの時系列的なデータが手に入ることになります。

あるいはまた，社会的なデータも時系列的に得られるものが多いです。一番わかりやすいのは株価の指標などでしょうか。日経平均株価など，時事刻々と変化する株価がグラフになって表されているのをみなさんもみたことがあると思います。また，TwitterなどのSNSで，今どのようなキーワードが使われているかといったことが，リアルタイムに計測されます。これらの指標は個々人の特徴というより，社会全体のうねりのようなものを表現していることになりますが，これを研究することも立派な社会心理学的テーマになり得ます。時系列データは，人文社会科学領域ではこれまで経済学やマーケティングなどが専門的に扱ってきたところがありますが，これからは心理学の領域でもこうしたデータを分析することが流行してくるかもしれません\footnote{もちろん生理心理学や動物心理学では，古くからこうしたデータを扱ってきています。}。

さて時系列的なデータはどういった特徴があるかと言うと，端的に言えば\keyterm{自己相関}{じこそうかん}{auto correlation}がある，ということでしょう。すなわち，ある時点のデータはその前の時点と相関している，ということです。体重のデータの場合でもそうですが，明日の体重は突然ランダムな値になるのではなく，今日の体重に応じて変化するはずだからです。また，データの変動に周期的なパターンが存在することもあります。脳波などのデータは基本的に波ですから，大きな波，小さな波がパターンを描いています。株価など社会的なデータも，季節による一定の変動など全体的な揺らぎがある上で，目的とする意味のある変化を見つけ出さなければならないという課題に取り組んでいます。たとえば周波数のデータの平均はゼロになりますから，基本的な記述統計では太刀打ちできないところがあります。回帰分析もその前提として，各データが独立に得られていると言うものがありますから，自己相関するデータは当然この仮定に違反していることになります。こうしたデータを分析するためには，周波数の大きさをスペクトル解析するとか，多次元の行列であるテンソルを使ってさまざまな要素を分解する，といった特殊な応用数学を駆使する必要があります。

今回はこうしたさまざまなモデルの中でも，\keyterm{状態空間モデル}{じょうたいくうかんもでる}{state space model}を紹介します。
\section{状態空間モデル}
時系列的なデータは，自己相関すなわち前の時点の状態が，次の時点に影響しているという特徴があり，これをうまく表現するのが状態空間モデルです。状態空間モデルはこれまでの\keytermL{モデリング}{もでりんぐ}の技法を使うと簡単に実装できます。というか，これまでの時系列的なデータ解析は，自己相関を減らすために一時点前のデータとの差分をとってそれをモデルにする，時点ごとの変化をスムージングするなど，さまざまな工夫の上に成り立っていました。状態空間モデルはそれを非常にすっきり表現してくれているので，それまでの分析ノウハウがなくとも誰でも簡単に利用できるモデルだといえるかもしれません。

状態空間モデルは，まず観測されたデータと，その背後にある\textbf{状態}を区別します。
体重変化のデータの例で考えてみましょう。計測された体重は，当然体の重さを反映したものですが，測定の時に多少の誤差が生じるでしょう。例に使っている筆者のデータでも，毎朝全裸で計測しているわけではありませんから，季節の変わり目で薄いパジャマからジャージにしたり，冬用の厚手のパジャマに変えたりすると，それだけで計測された数字には違いができます。本来知りたいのは，体の重さ，肉の重さそのものなはずですね。ここでいう体重そのもの，計測誤差のない肉の重さが，ここでいう状態のことになります。

そして$t$時点の状態$\mu_t$は，次の時点の状態に影響します。服の重さなどを取り除いた$t$時点の肉の状態に，食事や運動といった変化が加わって明日$t+1$時点の肉の状態になるわけですから，$\mu_{t+1} = f(\mu_t)$と考えることができるわけです。

もし翌日の体重が，今日の体重にちょっとした誤差がつくような変動しかしない，と言うのであれば$\mu_{t+1} \sim N(\mu_t,\tau)$と表すことができるでしょう。あるいは，運動量$P_t$の関数だというのであれば，$\mu_{t+1} \sim N(\mu_t + \beta_1P_t,\tau)$のように回帰分析の確率モデルのように表現すれば良いでしょう。もちろん説明変数が増えたり，一次関数ではないものを考えても構いません。

一方，計測値$W_t$は，この$\mu_t$に誤差$\sigma$がついて得られる，すなわち$W_t \sim N(\mu_t , \sigma)$という関係になります。これらの関係を表現したのが，図\ref{fig:29_01}になります。
\begin{figure}[ht]
	\begin{center}
		\includegraphics[width=0.8\linewidth,keepaspectratio]{../images/29_SSpace/slide29_01.png}
		\caption{状態空間モデルのイメージ(体重変化の例)}
		\label{fig:29_01}
	\end{center}
\end{figure}
このようにして，状態が時間$t,t+1,t+2,\cdots$を通じてつながっており，状態に応じて観測値が得られると考えます。観測値を手に入れる際には偶然誤差$\sigma$が生じますし，状態の変化は確率的に変わるのであれば幅$\tau$で，系統的に変わるのであればそれをモデリングしてやれば良いことになります。

この発想に基づいて，設計図を書いてみましょう。ここでは体重の変化に特別な傾向を考えず，ただ幅$\tau$で確率的にノイズが加わる\keyterm{ホワイトノイズモデル}{ほわいとのいずもでる}{white-noise model}で考えてみることにします。
少しイメージしにくいかもしれませんが，たとえば次のように表現できます(図\ref{fig:29_02})。
\begin{figure}[ht]
	\begin{center}
		\includegraphics[width=0.8\linewidth,keepaspectratio]{../images/29_SSpace/slide29_02.png}
		\caption{状態空間モデルの設計図}
		\label{fig:29_02}
	\end{center}
\end{figure}

それではこれをコードにしてみましょう。少し長くなりましたが，ホワイトノイズのモデルはコード\ref{stan::29_01}のようになります。
\begin{lstlisting}[language=C++,label={stan::29_01},caption={ホワイトノイズモデル}]
data{
 int L;
 array[L] real W;
}

parameters{
  real muZero;
  array[L] real mu;
  real<lower=0> sig;
  real<lower=0> tau;
}

model{
  mu[1] ~ normal(muZero, tau);
  
  for(l in 1:L){
    W[l] ~ normal(mu[l],sig);
  }
  
  for(i in 2:L){
    mu[i] ~ normal(mu[i-1],tau);
  }
  
  muZero ~ normal(80,10); 
  sig ~ cauchy(0,5);
  tau ~ cauchy(0,5);
}
\end{lstlisting}
\paragraph{コード解説}
\begin{description}
	\item[dataブロック] データ長L，各時点の体重Wがデータになります。
	\item[parametersブロック] データ長と同じだけの状態を表す\verb|mu|と，状態の変動幅\verb|tau|，状態に付加される計測誤差の変動幅\verb|sig|を宣言しています。それに加えて，最初の状態$\mu_0$をパラメータとしました。これは2時点目の状態$\mu_2$は$\mu_2 \sim N(\mu_1,\tau)$で，3時点目は$\mu_3 \sim N(\mu_2,\tau)$で・・・と表現できるのに対し，最初の点は$\mu_1 \sim N(\mu_0, \tau)$となって，計測される前の状態を考えなければならないからです。わからないものは確率で表現する，というベイズの流儀に則って，これはパラメータとして推定してやることにします。
	\item[modelブロック] まず1時点目の状態は，0時点目の状態から出てくるものですので，$\mu_1 \sim N(\mu_0,\tau)$をモデル化しました。次に計測された値$W_l$は状態$\mu_l$から出てくるものですから，1時点目からデータ長Lまで順に$W_l \sim N(\mu_l,\tau)$として尤度を書きます。次に状態のモデルですが，\underline{2時点目以降}は前の時点からの影響で表現できるので，$\mu_i \sim N(\mu_{i-1}, \tau)$としています。最後に事前分布として，これまでの経験から$\mu_0 \sim N(80,10)$とし，変化の幅はいずれもSDですから半コーシー分布で表現しました。
\end{description}

これを使って，体重のデータを分析してみましょう。データは2020年からのものを使います。図\ref{fig:29_03}をみると，2021年まではあまり変化がなく，21年初頭に少し体重が下がって，また上がって，その後5月から謎の下降\footnote{自転車事故による鎖骨骨折と，それに伴ってお箸持ちにくくて食事量が減る，運動も減るといったのが原因でしょう。}，そしてその反動\footnote{悔しいです。}，とおそらく線形モデルでは表現できないような複雑な動きをしています。

\begin{figure}[ht]
	\begin{center}
		\includegraphics[width=0.8\linewidth,keepaspectratio]{../images/29_SSpace/Rplot29_01.png}
		\caption{2020年の体重変化}
		\label{fig:29_03}
	\end{center}
\end{figure}


さて，コード\ref{code:29_01}を実行し，図\ref{fig:29_04}のような結果を得ます。複雑な動きについても，モデルがしっかりと予測しているのがみて取れますね。状態空間のパワフルな表現力をお楽しみいただけましたでしょうか。
\begin{lstlisting}[language=c++,label={code:29_01},caption={状態空間モデルのコード}]
dat1 <- dat %>%
  filter(date > "2020/01/01") %>%
  mutate(date = as.Date(date))

model <- cmdstan_model("StateSpace.stan")
dataSet <- list(L = NROW(dat1), W = dat1$weight)
fit1 <- model$sample(
  data = dataSet,
  chains = 4,
  parallel_chains = 4
)

fit1.stanfit <- fit1$output_files() %>% rstan::read_stan_csv()
fit1.df <- fit1.stanfit %>% as.data.frame() %>%
  as_tibble() %>%
  rowid_to_column("iter") %>%
  pivot_longer(-iter, names_to = "Varname") %>%
  group_by(Varname) %>%
  summarise(
    EAP = mean(value),
    MED = median(value),
    MAP = map_estimation(value),
    SD = sd(value),
    L95 = quantile(value, probs = 0.025),
    L50 = quantile(value, probs = 0.25),
    U50 = quantile(value, probs = 0.75),
    U95 = quantile(value, probs = 0.975)
  )

Est1 <- fit1.df %>%
  dplyr::filter(str_detect(Varname, "mu")) %>%
  dplyr::mutate(ID = str_extract(Varname, pattern = "\\d+") %>% 
  as.numeric()) %>%
  arrange(ID)

g <- dat1 %>%
  rowid_to_column("ID") %>%
  left_join(Est1, by = "ID") %>%
  ggplot(aes(x = ID, y = weight, ymin = U95, ymax = L95)) +
  geom_point() +
  geom_point(aes(x = ID, y = MAP), color = palette()[2]) +
  geom_ribbon(fill = palette()[3], alpha = 0.2)
plot(g)
\end{lstlisting}
\paragraph{コード解説}
\begin{description}
	\item[1-3行目] データファイルから該当する日程だけ抜き出します。日付を\verb|Date|形式に変換しています。
	\item[5-11行目] \verb|cmdstanr|でコンパイルしてサンプリングするところです。
	\item[13行目] \verb|cmdstanr|で作ったファイルを\verb|stanfit|オブジェクトに変換しています。\verb|rstan|パッケージを使っている人はこの行を実行する必要がありません。
	\item[14-28行目] MCMCサンプルを集計し，\keytermL{EAP}{EAP}，\keytermL{MED}{MED}，\keytermL{MAP}{MAP}推定値，SD，確信区間などを出しています。
	\item[30-34行目] 推定された状態$\mu_i$を取り出しています。
	\item[36-43行目] もとのデータと推定値を合体させ，プロットしています。
\end{description}

\begin{figure}[ht]
	\begin{center}
		\includegraphics[width=0.8\linewidth,keepaspectratio]{../images/29_SSpace/Rplot29_02.png}
		\caption{2020年の体重変化をモデルで追跡してみた結果}
		\label{fig:29_04}
	\end{center}
\end{figure}


\section{欠測値の補間}
ところで，データをよくみてみると，実は毎日のデータになっていないところがあることに気づきます。いや，気づかないかもしれないですね。次のコード\ref{code:29_02}を実行してみてください。
\begin{lstlisting}[language=c++,label={code:29_02},caption={日付は連続かな?}]
dat1 %>%
  mutate(lag = lag(date)) %>%
  mutate(date = as.Date(date), lag = as.Date(lag)) %>%
  mutate(FLG = date - lag) %>%
  dplyr::filter(FLG > 1)
\end{lstlisting}
\paragraph{コード解説}
\begin{description}
	\item[1行目] 先ほどの推定に使ったデータファイルです。
	\item[2行目] \verb|lag|関数で，データを1行ずらした列を作ります。
	\item[3行目] 体重を記録した日の変数\verb|date|と，先ほど作った一行ずらした変数\verb|lag|はいずれも日付に関する変数ですので，日付型に変更します。
	\item[4行目] 日付変数の引き算をして，一行下のデータと何日ずれていたのか算出しています。
	\item[5行目] 1日以上ずれている日をフィルタリングで抜き出しています。
\end{description}
これをみると，ちょくちょく抜けていて，時には4日も空いていたりすることがわかります\footnote{2020年2月22日から24日は岡山県にベイズ塾合宿(出張)でした。コロナがまだ豪華客船の中だけに抑え込まれている時期で，ギリギリ出張できた時期です。このあと大学の卒業式がなくなったり，前期オンライン授業になったりと言う，波乱の2020年度が始まったのでした。}。ともかく，これでは正確にデータを分析できているとは言えません。
\begin{Rscreen}{抜けている日}{out29_01}
	\begin{verbatim}
        # A tibble: 21 × 5
        date       weight bodyFat lag        FLG   
        <date>      <dbl>   <dbl> <date>     <drtn>
      1 2020-01-13   83.1    26.8 2020-01-11 2 days
      2 2020-02-01   82.6    26.7 2020-01-30 2 days
      3 2020-02-12   82.6    26.7 2020-02-10 2 days
      4 2020-02-26   83.1    26.8 2020-02-22 4 days
      5 2020-02-28   83.7    27   2020-02-26 2 days
      6 2020-03-02   83.2    26.9 2020-02-29 2 days
      7 2020-03-21   83.4    26.9 2020-03-19 2 days
      8 2020-05-09   82.8    26.8 2020-05-07 2 days
      9 2020-05-18   82.8    26.7 2020-05-16 2 days
     10 2020-07-31   82.8    27.7 2020-07-29 2 days
     # … with 11 more rows\end{verbatim}
\end{Rscreen}
しかし計測はできていなくても，出張中にも状態の変化は続きいているはずで，次の計測時点はその日の状態の関数になっているはずです(図\ref{fig:29_05})。
\begin{figure}[ht]
	\begin{center}
		\includegraphics[width=0.8\linewidth,keepaspectratio]{../images/29_SSpace/slide29_03.png}
		\caption{計測できない時点も変化は続く}
		\label{fig:29_05}
	\end{center}
\end{figure}

この欠損した状態のデータをどのように考えれば良いでしょうか。実はこのことのヒントはすでに，最初のコードの中に入っています。
状態空間モデルの1時点目，$\mu_1$は$\mu_0$から影響されているはずですが，データから$\mu_0$はわかりません。しかしわからないことは確率で表現するのがベイジアンのやり方です。
今回も，欠損している測定値はわからないものとして確率で表現し，推定してやれば良いのです！このようにして間を埋めることをとくに\keyterm{補間}{ほかん}{interpolating}と言います。

補間をするためには，Rコードの方でもStanコードの方でもすこし工夫が必要です。
まずRのほうで，データを加工するところから見てみましょう。Stanには欠測値\verb|NA|を与えることはできませんから，欠損しているところには特別なあり得ない数字，そうですね，999という数字でも入れておきましょう。

\begin{lstlisting}[language=c++,label={code:29_03},caption={連続したデータを作り欠測値に借りの値を代入する}]
fullDays <- 
  data.frame(date = as.Date("2020/01/01"):as.Date("2021/12/01")) %>%
  mutate(date = as.Date(date, origin = "1970-01-01")) %>%
  left_join(dat1, by = "date") %>%
  tidyr::replace_na(list(weight = 999, bodyFat = 999))
\end{lstlisting}
\paragraph{コード解説}
\begin{description}
	\item[2行目] 2020年1月1日から2021年12月1日までにデータを限定します。この期間，連続した日付を\verb|date|変数に追加します。
	\item[3行目] 日付変数を日付型に変更しています。
	\item[4行目] もとのデータを日付変数に結合しています。\verb|left_join|関数は左にもとのデータセットを置き，右から\verb|dat1|を引っ付けていきます。結合する際は，変数\verb|date|をキーにし，キー変数が同じものは同じ行に合わせるという湯やり方です。このやり方だと，該当する行がない場合(計測データセットの日付がない場合)，もとのデータセットは残して，体重や体脂肪変数は\verb|NA|になります。
	\item[5行目] \verb|replace_na|関数で，欠損のところに999という数字を入れています。
\end{description}

このようにすることで，次のようなデータセットができました(出力\ref{RS:out29_02})。
\begin{Rscreen}{抜けている日を無くした完全データセット}{out29_02}
	\begin{verbatim}
> fullDays
          date weight bodyFat
1   2020-01-01   83.3   26.70
2   2020-01-02   83.1   26.80
3   2020-01-03   83.2   26.90
4   2020-01-04   83.6   27.00
5   2020-01-05   82.9   28.70
6   2020-01-06   82.9   26.80
7   2020-01-07   82.7   26.70
8   2020-01-08   84.0   27.10
9   2020-01-09   84.6   27.30
10  2020-01-10   83.4   26.90
11  2020-01-11   84.3   27.20
12  2020-01-12  999.0  999.00
13  2020-01-13   83.1   26.80
14  2020-01-14   82.1   65.00
15  2020-01-15   82.7   26.70\end{verbatim}
\end{Rscreen}
これをみると，2020年1月12日は欠損だったのですが，データ上は999という数字が入って完全データのように見えます\footnote{ちなみに2020/01/11から12にかけては，犬会という研究会があったため関西に出張していました。宿に泊まったので，12日の朝のルーティンができなかったのです。}。このデータをStanに与えてやり，体重が999というあり得ない数字の場合は別の処理をする，という0過剰ポアソンの技術を応用します(セクション\ref{ZIpoisson}，Pp.\pageref{ZIpoisson})。
\begin{lstlisting}[language=c++,label={stan::29_02},caption={欠測値補間のコード}]
data{
 int L;
 array[L] real W;
 int<lower=0> Nmiss;
}

parameters{
  real muZero;
  array[L] real mu;
  array[Nmiss] real<lower=0> Miss_W;
  real<lower=0> sig;
  real<lower=0> tau;
}

model{
  mu[1] ~ normal(muZero, tau);
  
  {
    int j = 0;
    for(l in 1:L){
      if( W[l] != 999){
          // こっちは尤度
          W[l] ~ normal(mu[l],sig);
      }else{
        j = j + 1;
        // こっちはパラメータ
        Miss_W[j] ~ normal(mu[l], sig);
      }
    }
  }
  
  for(i in 2:L){
    mu[i] ~ normal(mu[i-1],tau);
  }
  
  muZero ~ normal(80,10); 
  sig ~ cauchy(0,5);
  tau ~ cauchy(0,5);
}
\end{lstlisting}

\begin{description}
	\item[dataブロック] データ長L，各時点の体重Wがデータです。加えて，推定すべき欠測値の数もデータとして受け取ります。
	\item[parametersブロック] データ長と同じだけの状態を表す\verb|mu|と，状態の変動幅\verb|tau|，状態に付加される計測誤差の変動幅\verb|sig|を宣言しています。また最初の状態$\mu_0$を表すパラメタと，欠損の数だけある推定用パラメータ\verb|Miss_W|を用意しました。
	\item[modelブロック] 先ほど同様，0時点目の状態から出てくるものですので，$\mu_1 \sim N(\mu_0,\tau)$をモデル化しています。次に各回のデータについて，$W_l$が999でなければ(同じでない，という論理式は\verb|!=|で表現します)，普通の状態空間モデルのように尤度として計算します。 そうではない，すなわちデータ$W_l=999$である，つまり欠測値であるはずということなら，$W_{miss} \sim N(\mu_l,\sigma)$で推定してやります。ここで特殊な変数$j$が出てきています。これは「何番目の欠測値か」を表す変数です。ブロックの中の一部だけで変数を宣言するため，まず\verb|for|文全体を中括弧で括り，はじめに$j=0$という数字を宣言しています。ここで欠測値推定のシーンになると，$j$のカウターを1つ増加させ，$j$番目の欠測値の推定をさせる，というやり方をしています。欠測値の数と行番号が合致しないので，こうした工夫が必要なのですね。その後のコードは先ほどと同じです。
\end{description}

このようにして，推定してみましょう。Stanに与えるデータセットは次のようにして作ります。
\begin{lstlisting}[language=c++,label={code:29_04},caption={Stanに与えるデータセット}]
dataSet <- list(L = NROW(fullDays), W = fullDays$weight, 
                Nmiss = sum(fullDays$weight == 999))
\end{lstlisting}
欠測値の数は，オブジェクト\verb|fullDays|の\verb|weight|変数が999かどうかを判定させ(\verb|==|でイコールかどうかという論理判断になります)，条件が合致した数を数え上げるという方法をとっています。

コードが走り出すと，推定は順調に進むと思います。が，結果の出方が少しややこしいです。状態$\mu_i$はすべての日程について推定されますが，これに加えて欠損があったデータも推定されていきますので，描画する際は「実測値なのか推定値なのか」，「補間した欠損の推定値は何番目の欠測値だったか」を判断しながら結合していき，データにする必要があるからです。Stanの中で作ったカウント変数$j$をRの方でも考えてやらないといけないわけですね。
\begin{lstlisting}[language=c++,label={code:29_04b},caption={プロット用にデータを整形する関数とプロットのコード}]
Est2 <- fit2.df %>%
  dplyr::filter(str_detect(Varname, "mu")) %>%
  dplyr::filter(!str_detect(Varname, "muZero")) %>%
  dplyr::mutate(ID = str_extract(Varname, pattern = "\\d+") %>% 
  as.numeric()) %>%
  arrange(ID) %>%
  select(ID, MAP, U95, L95)

Est2miss <- fit2.df %>%
  dplyr::filter(str_detect(Varname, "Miss_W")) %>%
  dplyr::mutate(ID = str_extract(Varname, pattern = "\\d+") %>% 
  as.numeric()) %>%
  arrange(ID) %>%
  select(ID, MAP, U95, L95)

### plot用の関数を準備
plotFunction <- function(fullDays, Est, MissEst) {
  tmp <- fullDays %>%
    rowid_to_column("ID") %>%
    left_join(Est, by = "ID") %>%
    rowwise() %>%
    mutate(FLG = if (weight != 999) {1} else {2})
  misJ <- 1
  tmp$weight2 <- NA
  tmp$weight2U <- NA
  tmp$weight2L <- NA
  for (i in 1:NROW(tmp)) {
    if (tmp$FLG[i] == 2) {
      tmp$weight2[i] <- MissEst$MAP[misJ]
      tmp$weight2U[i] <- MissEst$U95[misJ]
      tmp$weight2L[i] <- MissEst$L95[misJ]
      misJ <- misJ + 1
    } else {
      tmp$weight2[i] <- tmp$weight[i]
      tmp$weight2U[i] <- tmp$weight[i]
      tmp$weight2L[i] <- tmp$weight[i]
    }
  }
  return(tmp)
}

plot.tmp <- plotFunction(fullDays, Est2, Est2miss)
g <- ggplot(data = plot.tmp) +
  geom_point(aes(x = date, y = weight2)) +
  geom_errorbar(aes(x = date, y = weight2, ymin = weight2L, 
                    ymax = weight2U, color = palette()[2])) +
  geom_point(aes(x = date, y = MAP, color = palette()[3])) +
  geom_errorbar(aes(x = date, y = MAP, ymin = L95, 
                    ymax = U95, color = palette()[4])) +
  scale_x_date(date_breaks = "1 month", 
                limits = as.Date(c("2020-01-01", "2020-05-01"))) +
  theme(legend.position = "none")
\end{lstlisting}
\paragraph{プログラム解説}
\begin{description}
	\item[Est2を作るブロック] MCMCサンプルをデータフレーム化したものの中から，状態$\mu$に関するデータだけを抜き出したものを作る。
	\item[Est2missを作るブロック] 欠測値を補間するための推定値を，MCMCサンプルから取り出した欠測値補間だけのデータセットを作る。
	\item[plot用の関数] プロットはこの後にも行いますので，もう関数を作ってまとめてしまうことにします。もとのデータセット\verb|fullDays|と，状態の推定値データセット\verb|Est|，欠測値の推定値データセット\verb|MissEst|を引数にとる関数です。
	      \begin{description}
		      \item[関数3行目] まずもとデータに行番号を\verb|ID|として変数化します。
		      \item[関数4行目] 状態の推定値は，すべての行について推定しているはずですので，行番号をキーに結合\verb|left_join|します。
		      \item[関数5行目] もとの観測値が実測値なのか欠測値なのかを表現する変数\verb|FLG|を用意しておきます。これらはすべて，一時的なオブジェクト\verb|tmp|に納められます。
		      \item[関数6行目] 欠測値のインデックスを表す変数\verb|misJ|をつくり，カウントを1に設定しておきます。
		      \item[関数7-9行目] 先ほどの一次オブジェクトにいまから体重の推定値を入れていきますが，この推定値は実測値で得られている時は不要で\verb|NA|になるはずですので，いったんすべてに\verb|NA|を代入しています。
		      \item[関数10-21行目] オブジェクト\verb|tmp|を一行ずつ見ていって，もし\verb|FLG|変数が欠測値であることを教えてくれたら，\verb|misJ|番目の推定値を代入します。代入が終わったらカウンタを1つ追加しておきます。欠測値でなければ，もとの体重をそのまま移し入れます。ちなみにUやLがついているのは95\%区間の上限と下限なのですが，実測値の場合は推定ではありませんのですべて同じ数字になります。
		      \item[関数22行目] 作ったオブジェクト\verb|tmp|を戻り値として返します。
	      \end{description}
	\item[関数の結果を受けとる] \verb|plot.tmp|は描画用の一時データオブジェクトです。ここに全期間のデータと，今回の状態推定値，欠損推定値を入れておきます。
	\item[描画する] x軸は全日程です。まずは体重$W_j$を\verb|geom_point|と\verb|geom_errorbar|でプロットします。欠損の場合は推定値が入っています。次のポイントとエラーバーのgeomは，状態$\mu$のものになります。
\end{description}

さあどうでしょう。コードはともかく，図\ref{fig:29_05}を見てみましょう。実際に観測しているのは黒点，状態の推定値は緑点と赤いエラーバーで表示されています。青いバーは欠損した時の推定値で，値がないのですから当然触れ幅は大きいですが，おそらくこの体重が測定できなかった日はこれぐらいだったんじゃないか，という値が推定されているわけです。数字で見ると，たとえば2020年1月12日のMAP推定値は83.4kg，95\%区間でいうなら[82.6,84.2]です。前日が84.3kgで翌日が83.1kgですから，そんなもんかなあという気もします。
\begin{figure}[ht]
	\begin{center}
		\includegraphics[width=0.8\linewidth,keepaspectratio]{../images/29_SSpace/Rplot29_03.png}
		\caption{欠損していても推定できる}
		\label{fig:29_06}
	\end{center}
\end{figure}

ところで，測定していないはずの日も補間できるよ，ということを考えてみると，これはすごいことだと思いませんか。そうです，未来の値ももちろん「測定していない日」ですから，このモデルを使うと未来予測ができるのです！すごい！

早速同じコードで，日程を未来に飛ばしてみましょう。データは2021年12月06日までしかありません\footnote{この原稿は2021年12月10日17時に執筆しています。}から，2021年12月31日までデータを伸ばし，推定してみようではありませんか。
コードは省略しますが，実行した結果は次のような図\ref{fig:29_07}として得られます。
\begin{figure}[ht]
	\begin{center}
		\includegraphics[width=0.8\linewidth,keepaspectratio]{../images/29_SSpace/Rplot29_04.png}
		\caption{未来もある意味欠測値}
		\label{fig:29_07}
	\end{center}
\end{figure}

これをみると，データが得られなくなった次の日から，確信区間が日を追うごとに広がっていく様がみて取れます。当然そうですね。$\mu_{t+1}$は$\mu_t$にプラスかマイナスか，いずれとも言えない振幅幅で揺れ動くのですから，先に行けば行くほどその可能性は不確かなものになっていくのです。ちなみに描画用のデータから2021年最後の数日間の予測値を見てみると，次のようになっています。
\begin{lstlisting}[language=c++,label={code:29_05},caption={予想される2021年末の体重(一部略)}]
> plot.tmp %>% tail
# A tibble: 6 × 11
# Rowwise: 
   ID date       weight bodyFat   MAP   U95   L95   FLG weight2
  <dbl> <date>    <dbl>   <dbl> <dbl> <dbl> <dbl> <dbl>   <dbl>
  738 2021-12-26    999      NA  82.2  83.8  80.5     2    82.1
  739 2021-12-27    999      NA  82.3  83.9  80.4     2    82.1
  740 2021-12-28    999      NA  82.2  83.9  80.4     2    82.2
  741 2021-12-29    999      NA  82.2  84.0  80.4     2    82.0
  742 2021-12-30    999      NA  82.3  84.0  80.3     2    82.3
  743 2021-12-31    999      NA  82.3  84.0  80.3     2    82.2d
\end{lstlisting}
未来の状態$\mu$は82.2のあたりでほぼ固定，それに伴って想定される体重\verb|weight2|も82.2程度ですが，幅がどんどんと広くなっていってます。可能性として，80kgに落ちているかもしれませんし，84kgになっているかもしれない，というのが現時点での予測になります。

予測ができるとは言え，なんだかほとんど変わり映えしない数字で面白くないですね。それは当然，「明日の体重は今日の体重+誤差だろう」という，とくに何らメカニズムや仮定を入れなかったものですから，ほぼ横一線の推定にしかならないわけです。

では少し欲張って，構造を入れみましょうか。
ここまでは，$\mu_t$は$\mu_{t-1}$にのみ影響される，というモデルになっていましたが，もう1日前の影響が入っていることを考えます。
たとえば体重が増加傾向にあるとか，減少傾向にあるといった，変化を考えてみることにしましょう。すなわち，$\mu_t - \mu_{t-1} = \mu_{t-1}-\mu_{t-2}+\varepsilon$です。これは昨日から今日への変化($\mu_t - \mu_{t-1}$)は，一昨日から昨日への変化($\mu_{t-1}-\mu_{t-2}$)と同じようなものだろう(誤差$\varepsilon$はあるけど)，と考えていることになります。このようなモデルを\keytermL{2階差分のトレンド}{2かいさぶんのとれんど}と呼びます。

この式を変形すると，
\[ \mu_t = \mu_{t-1} + \mu_{t-1} -\mu_{t-2} + \varepsilon \]
\[  = 2\mu_{t-1} -\mu_{t-2} + \varepsilon \]
となることはすぐにわかりますね。これを確率モデルで表現するなら，
\[ \mu_t \sim N(2\mu_{t-1}-\mu_{t-2},\tau) \]

と考えていることと同じですから，そのようにStanのコードも変形します(コード\ref{stan::29_03})。

\begin{lstlisting}[language=c++,label={stan::29_03},caption={2階差分トレンドのコード}]
...(前略)...
model{
  mu[1] ~ normal(muZero, tau);
  mu[2] ~ normal(mu[1] , tau);
  
  {
    int j = 0;
    for(l in 1:L){
      if( W[l] != 999){
          // こっちは尤度
          W[l] ~ normal(mu[l],sig);
      }else{
        j = j + 1;
        // こっちはパラメータ
        Miss_W[j] ~ normal(2*mu[l-1]-mu[l-2], sig);
      }
    }
  }
  
  for(i in 3:L){
    //2階差分
    mu[i] ~ normal(2*mu[i-1]-mu[i-2],tau);
  }
  
  muZero ~ normal(80,10); 
  sig ~ cauchy(0,5);
  tau ~ cauchy(0,5);
}
\end{lstlisting}
コードは2階差分になっているので，状態は3時点目からしか推定できません。最初の2行で$\mu_1,\mu_2$をメカニズムに沿って推定し，変化が取れるようになってから，は2階差分のコードで推定します。\verb|for|文が3から始まっていることに注意してください。

このコードと年末までのデータ，そしてプロット関数を使って推定した結果が次の図\ref{fig:29_08}になります。
\begin{figure}[ht]
	\begin{center}
		\includegraphics[width=0.8\linewidth,keepaspectratio]{../images/29_SSpace/Rplot29_05.png}
		\caption{2階差分トレンドの予測}
		\label{fig:29_08}
	\end{center}
\end{figure}

これを見ると，それまでの増加傾向を反映して，年末には84kg，最悪の場合86.6kgまで増えている可能性があることがわかります\footnote{なんてこった，まだ執筆中は答えが出ていませんが，この予測通りになったらとんでもないことですよ。だって私は20代66kgだったんです。あの頃から20kgも増えている自分の未来なんて想像したくないですよ。またダイエット始めるかなぁ。}\footnote{2021/01/15追記。2021/12/31,6:07AMの記録は82.3kg,体脂肪率26.6\%でした。ちゃんと確信区間の中に入ってますね！}。

\begin{lstlisting}[language=c++,label={code:29_06},caption={2階差分トレンドで予想される2021年末の体重(一部略)}]
> plot.tmp %>% tail()
# A tibble: 6 × 11
# Rowwise: 
   ID date       weight bodyFat   MAP   U95   L95   FLG weight2
<dbl> <date>    <dbl>   <dbl> <dbl> <dbl> <dbl> <dbl>   <dbl>
  738 2021-12-26    999      NA  83.9  85.6  81.5     2    83.6
  739 2021-12-27    999      NA  83.9  85.8  81.5     2    83.8
  740 2021-12-28    999      NA  84.0  85.9  81.4     2    83.8
  741 2021-12-29    999      NA  84.1  86.1  81.4     2    83.9
  742 2021-12-30    999      NA  84.2  86.3  81.3     2    84.1
  743 2021-12-31    999      NA  84.2  86.5  81.2     2    84.0
\end{lstlisting}

今回は2階差分のトレンドを入れての予測となりましたが，この他にも周期的に影響してくる季節項をいれるとか，たとえば「週末は体重が増加するだろう」と考えて週末効果の項を入れるとか，さまざまな工夫を思いつくことができるのではないでしょうか。

\section{状態空間モデルの展開}
最後になりましたが，状態空間モデルは表面に現れている値の背後に潜在変数を仮定し，それを滑らかに繋ぎ合わせた\keyterm{スムージング}{すむーじんぐ}{smoothing}の技術であるとも考えることができます。時系列を通じて，その前後の「状態」を，柔らかなバネで結合し，データ全体を通じて変化する関数として相互に調整させあっている，と考えることもできるでしょう。

また，時間というのは過去から未来へ進む，1次元の道のりです。ところで私たちは2次元の地図を見たり三次元の空間を見たりしていますよね。たとえばこの状態空間が，時系列的一次元ではなく，X軸とY軸に広がる網のような二次元結合していることを考えると，空間的な統計分析が可能になります。たとえば地価というのは，駅からの距離や大型商業施設，コンビニエンスストアなどの利便性に応じて変化しますが，おそらくそれらは連綿と繋がっていて，徐々に変化していくものであるはずです。であれば徐々に変化する要素をつなぎ合わせて，面の状態空間モデルを作ってやることもできるわけです。さらにその地価の変動を考える，つまり二次元平面$\times$時系列に潜在変数をつなげてスムージングしたモデルを作れば，時系列的な地価の変動を考えたり予測したりできるかもしれません。

心理学はこれまで，調査法で一時点を切り出してその構造を把握する，というアプローチが多くありました。時系列を扱うとしても，プレ・ポストの変化を見るとか，あまり長期にわたって追いかけるようなことはしてきませんでした。しかし冒頭でお話ししたように，これからの心理学データは時間的にも空間的にも広がるものを計測し，利用できるようになるかもしれません。そのためにはより深い人間の観察や推察によるモデルの作り込みと，状態空間モデルとベイズ推定のような数理モデルと強力な推定法のタッグが，必要になってくるでしょう。

\section{課題}
体重変化のデータには，体脂肪の記録もあります。体重$\times$体脂肪率から，筆者の筋肉量を計算できます。
2020年以降のデータを使って，筋肉量の変化を状態空間モデルで推定してみてください。2階差分トレンドモデルで，未来のデータも予測できればなお結構です。
これらのモデルを分析するR/Stanコードを提出してください。結果の解釈などを，スクリプトのコメントアウトや別添ファイルなどで提供してもらえるとなお良いです。もちろんRmdファイルでの提出であれば完璧です。なお提出されたコード単体でバグがなく動くことが確認できないものは，未提出扱いになります。コードの書き方などわからないところがあれば，曜日別TAか小杉までメールで連絡し，指導を受けてください。

\end{document}
